\documentclass[10pt,a4paper]{amsart}
\usepackage{amsmath,amssymb,amsthm,graphicx}
% Retained minimal preamble for the Stacks Project chapter more-algebra.tex (More on Algebra).
% The chapter file pinned-original-source.tex carries only \input{preamble} on its first line, so no
% preamble text lives inside the chapter itself.  The Stacks Project uses one single shared preamble for
% every chapter; the snapshot stages it identically as sources/stacks-chow/preamble.tex and
% sources/stacks-etale-cohomology/preamble.tex.  The block below is copied mechanically, line for line,
% from sources/stacks-chow/preamble.tex lines 14-18 (the Xy-pic setup) and 154-217 (the native theorem
% declarations together with the mathematical macros).  Nothing was added, removed or redefined; the
% wrapper class and the amsmath/amssymb/amsthm/graphicx base packages are supplied by the assembler.
% Chapter: more-algebra.tex -- More on Algebra.
% For commutative diagrams we use Xy-pic
\usepackage[all]{xy}

% We use 2cell for 2-commutative diagrams.
\xyoption{2cell}

% Theorem environments.
%
\theoremstyle{plain}
\newtheorem{theorem}[subsection]{Theorem}
\newtheorem{proposition}[subsection]{Proposition}
\newtheorem{lemma}[subsection]{Lemma}

\theoremstyle{definition}
\newtheorem{definition}[subsection]{Definition}
\newtheorem{example}[subsection]{Example}
\newtheorem{exercise}[subsection]{Exercise}
\newtheorem{situation}[subsection]{Situation}

\theoremstyle{remark}
\newtheorem{remark}[subsection]{Remark}
\newtheorem{remarks}[subsection]{Remarks}

\numberwithin{equation}{subsection}

% Macros
%
\def\lim{\mathop{\mathrm{lim}}\nolimits}
\def\colim{\mathop{\mathrm{colim}}\nolimits}
\def\Spec{\mathop{\mathrm{Spec}}}
\def\Hom{\mathop{\mathrm{Hom}}\nolimits}
\def\Ext{\mathop{\mathrm{Ext}}\nolimits}
\def\SheafHom{\mathop{\mathcal{H}\!\mathit{om}}\nolimits}
\def\SheafExt{\mathop{\mathcal{E}\!\mathit{xt}}\nolimits}
\def\Sch{\mathit{Sch}}
\def\Mor{\mathop{\mathrm{Mor}}\nolimits}
\def\Ob{\mathop{\mathrm{Ob}}\nolimits}
\def\Sh{\mathop{\mathit{Sh}}\nolimits}
\def\NL{\mathop{N\!L}\nolimits}
\def\CH{\mathop{\mathrm{CH}}\nolimits}
\def\proetale{{pro\text{-}\acute{e}tale}}
\def\etale{{\acute{e}tale}}
\def\QCoh{\mathit{QCoh}}
\def\Ker{\mathop{\mathrm{Ker}}}
\def\Im{\mathop{\mathrm{Im}}}
\def\Coker{\mathop{\mathrm{Coker}}}
\def\Coim{\mathop{\mathrm{Coim}}}

% Boxtimes
%
\DeclareMathSymbol{\boxtimes}{\mathbin}{AMSa}{"02}

%
% Macros for moduli stacks/spaces
%
\def\QCohstack{\mathcal{QC}\!\mathit{oh}}
\def\Cohstack{\mathcal{C}\!\mathit{oh}}
\def\Spacesstack{\mathcal{S}\!\mathit{paces}}
\def\Quotfunctor{\mathrm{Quot}}
\def\Hilbfunctor{\mathrm{Hilb}}
\def\Curvesstack{\mathcal{C}\!\mathit{urves}}
\def\Polarizedstack{\mathcal{P}\!\mathit{olarized}}
\def\Complexesstack{\mathcal{C}\!\mathit{omplexes}}
% \Pic is the operator that assigns to X its picard group, usage \Pic(X)
% \Picardstack_{X/B} denotes the Picard stack of X over B
% \Picardfunctor_{X/B} denotes the Picard functor of X over B
\def\Pic{\mathop{\mathrm{Pic}}\nolimits}
\def\Picardstack{\mathcal{P}\!\mathit{ic}}
\def\Picardfunctor{\mathrm{Pic}}
\def\Deformationcategory{\mathcal{D}\!\mathit{ef}}
\usepackage{cleveref}
\crefname{originallemma}{original lemma}{original lemmas}
\Crefname{originallemma}{Original lemma}{Original lemmas}
\makeatletter
\newlabel{derived-lemma-morphisms-from-projective-complex}{{064B}{source}{}{sourcepointer-1}{}}
\newlabel{derived-lemma-morphisms-from-projective-complex@cref}{{[originallemma][0][]064B}{[1][0][]source}{}{}{}}
\makeatother
\crefname{originallemma}{original lemma}{original lemmas}
\Crefname{originallemma}{Original lemma}{Original lemmas}
\makeatletter
\newlabel{derived-lemma-morphisms-into-injective-complex}{{05TG}{source}{}{sourcepointer-2}{}}
\newlabel{derived-lemma-morphisms-into-injective-complex@cref}{{[originallemma][0][]05TG}{[1][0][]source}{}{}{}}
\makeatother
\title{Maps out of an almost free module to a bounded below complex in the derived category}
\author{The Stacks Project Authors}
\date{stacks-project:more-algebra.tex:local-snapshot-20261012; GNU FDL 1.2 or later}
\begin{document}
\maketitle
\setcounter{subsection}{1}
\noindent\textit{Editorial note.} All delivered mathematical text is the author's own, copied from the pinned Stacks Project chapter file more-algebra.tex; no statement or proof step was rewritten, reordered or removed.\par\medskip
\noindent\textit{Editorial note.} The chapter file pinned-original-source.tex carries only \textbackslash{}input\{preamble\} on its first line, so no preamble text lives inside the chapter. The retained file sources/stacks-more-algebra/retained-preamble.tex is a mechanical, line-for-line copy of the project's single shared preamble as staged in the snapshot; nothing was added, removed or redefined and no package was added.\par\medskip
\noindent\textit{Editorial note.} The source URL is the source-level Stacks Project URL, as required when more than one card is produced from one source; the per-claim tag is carried by claim\_key.\par\medskip
\noindent\textit{Editorial note.} The two cross-chapter references inside the counted proof span (derived-lemma-morphisms-from-projective-complex and derived-lemma-morphisms-into-injective-complex) are recorded as non-load-bearing navigation pointers; the argument that consumes them stays inside the counted proof span.\par\medskip
\noindent\textit{Allowed standard prerequisites (editorial).}\begin{itemize}
\item Almost free modules and their filtrations by free modules. fixed earlier in the same section and used throughout the counted spans. Supplies the filtration along which the lift is constructed.
\item Bounded below complexes and the standard representatives of their derived objects. the bounded below hypothesis in the counted statement span, and the two cited derived-category lemmas in the counted proof span. Is what makes the successive obstructions computed in the lift vanish.
\end{itemize}\medskip
\ifdefined\hypertarget\hypertarget{sourcepointer-1}{}\fi
\noindent\textit{Source reference.} Original lemma 064B: focused-seeds/AGP/sources/stacks-snapshot-20261011/tags.txt, source lines 7860--7860. Cross-chapter navigation pointer into the Stacks Project; the standard fact it supplies is applied inside the counted proof span and the argument resting on it stays there.\par\medskip
\ifdefined\hypertarget\hypertarget{sourcepointer-2}{}\fi
\noindent\textit{Source reference.} Original lemma 05TG: focused-seeds/AGP/sources/stacks-snapshot-20261011/tags.txt, source lines 7480--7480. Cross-chapter navigation pointer into the Stacks Project; the standard fact it supplies is applied inside the counted proof span and the argument resting on it stays there.\par\medskip
\begin{lemma}
\label{lemma-map-out-of-almost-free}
Let $R$ be a ring. Let $M^\bullet$ be a complex of modules over $R$
with $M^i = 0$ for $i > 0$ and $M^0$ a projective $R$-module.
Let $K^\bullet$ be a second complex.
\begin{enumerate}
\item Assume $K^i = 0$ for $i \leq -2$. Then
$\Hom_{D(R)}(M^\bullet, K^\bullet) = \Hom_{K(R)}(M^\bullet, K^\bullet)$.
\item Assume $K^i = 0$ for $i \not \in [-1, 0]$ and
$K^0$ a projective $R$-module. Then for a map of complexes
$a^\bullet : M^\bullet \to K^\bullet$, the following are equivalent
\begin{enumerate}
\item $a^\bullet$ induces the zero map $\Ext^1_R(K^\bullet, N) \to
\Ext^1_R(M^\bullet, N)$ for all $R$-modules $N$, and
\item there is a map $h^0 : M^0 \to K^{-1}$ such that
$a^{-1} + h^0 \circ d^{-1}_K = 0$.
\end{enumerate}
\item Assume $K^i = 0$ for $i \leq -3$. Let
$\alpha \in \Hom_{D(R)}(M^\bullet, K^\bullet)$. If the
composition of $\alpha$ with
$K^\bullet \to K^{-2}[2]$ comes from an $R$-module map
$a : M^{-2} \to K^{-2}$ with $a \circ d_M^{-3} = 0$, then
$\alpha$ can be represented by a map of complexes
$a^\bullet : M^\bullet \to K^\bullet$ with $a^{-2} = a$.
\item In (2) for any second map of complexes
$(a')^\bullet : M^\bullet \to K^\bullet$
representing $\alpha$ with $a = (a')^{-2}$
there exist $h^i : M^i \to K^{i - 1}$ for $i = 0, -1$ such that
$$
h^{-1} \circ d_M^{-2} = 0, \quad
(a')^{-1} = a^{-1} + d_K^{-2} \circ h^{-1} + h^0 \circ d_M^{-1},\quad
(a')^0 = a^0 + d_K^{-1} \circ h^0
$$
\end{enumerate}
\end{lemma}

\begin{proof}
Set $F^0 = M^0$.
Choose a free $R$-module $F^{-1}$ and a surjection $F^{-1} \to M^{-1}$.
Choose a free $R$-module $F^{-2}$ and a surjection
$F^{-2} \to M^{-2} \times_{M^{-1}} F^{-1}$. Continuing in this
way we obtain a quasi-isomorphism $p^\bullet : F^\bullet \to M^\bullet$
which is termwise surjective and with $F^i$ projective for all $i$.

\medskip\noindent
Proof of (1). By
Derived Categories, Lemma \ref{derived-lemma-morphisms-from-projective-complex}
we have
$$
\Hom_{D(R)}(M^\bullet, K^\bullet) = \Hom_{K(R)}(F^\bullet, K^\bullet)
$$
If $K^i = 0$ for $i \leq -2$, then any morphism of complexes
$F^\bullet \to K^\bullet$ factors through $p^\bullet$. Similarly, any
homotopy $\{h^i : F^i \to K^{i - 1}\}$ factors through $p^\bullet$.
Thus (1) holds.

\medskip\noindent
Proof of (2). If (2)(b) holds, then $a^\bullet$ is homotopic to a map
of complexes $(a')^\bullet : M^\bullet \to K^\bullet$
which is zero in degree $-1$.
On the other hand, let $N \to I^\bullet$ be an injective resolution.
We have
$$
\Ext^1_R(K^\bullet, N) = \Hom_{D(R)}(K^\bullet, I^\bullet[1]) =
\Hom_{K(R)}(K^\bullet, I^\bullet[1])
$$
by Derived Categories, Lemma
\ref{derived-lemma-morphisms-into-injective-complex}.
Let $b^\bullet : K^\bullet \to I^\bullet[1]$ be a map of complexes.
Since $K^1 = 0$ the map $b^0 : K^0 \to I^1$ maps into the kernel of
$I^1 \to I^2$ which is the image of $I^0 \to I^1$. Since $K^0$ is projective
we can lift $b^0$ to a map $h : K^0 \to I^0$. Thus we see that $b^\bullet$
is homotopic to a map of complexes $(b')^\bullet$ with $(b')^0 = 0$. Since
$K^i = 0$ for $i \not \in [-1, 0]$ it follows that
$(b')^\bullet \circ (a')^\bullet = 0$ as a map of complexes.
Hence the map $\Ext^1_R(K^\bullet, N) \to
\Ext^1_R(M^\bullet, N)$ is zero. In this way we see that (2)(b)
implies (2)(a). Conversely, assume (2)(a). We see that the
canonical element in $\Ext^1_R(K^\bullet, K^{-1})$ maps to
zero in $\Ext^1_R(M^\bullet, K^{-1})$. Using (1) we
see immediately that we get a map $h^0$ as in (2)(b).

\medskip\noindent
Proof of (3). Choose $b^\bullet : F^\bullet \to K^\bullet$ representing
$\alpha$. The composition of $\alpha$ with $K^\bullet \to K^{-2}[2]$ is
represented by $b^{-2} : F^{-2} \to K^{-2}$. As this is homotopic to
$a \circ p^{-2} : F^{-2} \to M^{-2} \to K^{-2}$, there is a map
$h : F^{-1} \to K^{-2}$ such that $b^{-2} = a \circ p^{-2} + h \circ d_F^{-2}$.
Adjusting $b^\bullet$ by $h$ viewed as a homotopy from $F^\bullet$
to $K^\bullet$, we find that $b^{-2} = a \circ p^{-2}$. Hence $b^{-2}$
factors through $p^{-2}$. Since $F^0 = M^0$ the kernel of $p^{-2}$
surjects onto the kernel of $p^{-1}$ (for example because the kernel
of $p^\bullet$ is an acyclic complex or by a diagram chase). Hence $b^{-1}$
necessarily factors through $p^{-1}$ as well and we see that (3)
holds for these factorizations and $a^0 = b^0$.

\medskip\noindent
Proof of (4) is omitted. Hint: There is a homotopy between
$a^\bullet \circ p^\bullet$ and $(a')^\bullet \circ p^\bullet$
and we argue as before that this homotopy factors through $p^\bullet$.
\end{proof}

\end{document}
